\documentclass[11pt]{article}
\usepackage{amsmath,amssymb,amsthm,mathtools}
\usepackage[margin=1.1in]{geometry}
\usepackage{booktabs}
\usepackage{hyperref}

\theoremstyle{plain}
\newtheorem{theorem}{Theorem}
\newtheorem{proposition}[theorem]{Proposition}
\newtheorem{corollary}[theorem]{Corollary}
\newtheorem{lemma}[theorem]{Lemma}
\theoremstyle{definition}
\newtheorem{definition}[theorem]{Definition}
\newtheorem{remark}[theorem]{Remark}
\newtheorem{question}[theorem]{Question}

\newcommand{\Gam}{\Gamma}
\newcommand{\wo}{w_0}
\newcommand{\Sn}{\mathfrak{S}_n}
\DeclareMathOperator{\type}{type}

\title{The fibred bijection of Lenart--Sottile's Corollary 4 exists,\\
and its existence is empty}
\author{Clio}
\date{6 October 2026}

\begin{document}
\maketitle

\begin{abstract}
Lenart and Sottile (\texttt{math/0202090}, Corollary 4 and the remark following
it) observe that
$I_\alpha(u,w)=\sum_v c^w_{u,v}\,I_\alpha(\wo v,\wo)$ and write that a
type-preserving map $\Gam(u,w)\to\bigsqcup_v\Gam(\wo v,\wo)$ with fibre
cardinality $c^w_{u,v}$ ``would give a combinatorial interpretation for
$c^w_{u,v}$, and thus solve the Littlewood--Richardson problem.''  We prove that
such a map always exists, that the number of them is an explicit multinomial
product, and hence that \emph{the existence of the map is logically equivalent to
the counting identity it is derived from} and gives no interpretation of
$c^w_{u,v}$ whatsoever (Theorem~\ref{thm:A}).  We prove further that the
fibre-cardinality condition is vacuous for $n\le 5$, because every Schubert
structure constant of $\mathrm{Fl}_n$ is $0$ or $1$ for $n\le5$, with $c=2$ first
occurring at $n=6$ (Theorem~\ref{thm:B}).  We record that the
bijectivization-uniqueness clause of Petrov (\texttt{2609.18502}, \S5) does not
apply (Proposition~\ref{prop:C}), and that the natural candidate map --- the
identity on second-coordinate words --- is refuted.  The upshot is a correction
of the question: what is open is not existence but \emph{canonicity}, and the
smallest instances with any fibre content live in $\mathfrak{S}_6$.
\end{abstract}

\section{The statement, and the conventions it is stated in}

\subsection{Source}

\texttt{math/0202090} (Lenart--Sottile, \emph{Skew Schubert polynomials}) is held
at \texttt{deep-read}.  The theorem numbers were resolved from the compiled
\texttt{.aux} of \texttt{skewschub.tex} (the \LaTeX{} source carries a single
counter shared across theorem/proposition/lemma/corollary/example/remark, so
numbers \emph{must not} be counted by hand): \texttt{cor:identity} $=\{4\}\{4\}$,
i.e.\ Corollary~4 on p.~4, and \texttt{thm:skew} $=\{2\}\{4\}$.  The quoted
remark follows Theorem~2 at l.~486--502.

\subsection{Definitions}

Let $\Sn$ be the symmetric group, $\ell$ the Coxeter length, $\wo$ the longest
element, and $t_{ij}$ the transposition of $i<j$.  Products are composition,
$(pq)(i)=p(q(i))$.

\begin{definition}[labelled Bruhat order; \texttt{math/0202090} l.~352]
For each Bruhat cover $u\lessdot w$ with $u^{-1}w=t_{ij}$, $i<j$, put $j-i$
directed edges $u\xrightarrow{(k,b)}w$, one for each integer $k$ with
$i\le k<j$, where $b=u(i)=w(j)$.
\end{definition}

\begin{definition}[increasing chain, type; l.~366, l.~372]
An \emph{increasing chain} $\gamma$ from $u$ to $w$ is a saturated chain
$u=u_0\lessdot u_1\lessdot\dots\lessdot u_m=w$ together with a choice of label on
each step such that the label sequence $(k_1,b_1),\dots,(k_m,b_m)$ is strictly
increasing in lexicographic order on pairs.  The \emph{type} of $\gamma$ is the
exponent vector $\alpha=\type(\gamma)\in\mathbb{Z}_{\ge0}^{n-1}$ with
$\alpha_i=\#\{s: k_s=i\}$.  Write $\Gam(u,w)$ for the set of increasing chains
from $u$ to $w$, $\Gam_\alpha(u,w)$ for its type-$\alpha$ part, and
$I_\alpha(u,w)=|\Gam_\alpha(u,w)|$.  Note $|\alpha|=m=\ell(w)-\ell(u)$.
\end{definition}

Let $c^w_{u,v}$ be the Schubert structure constants of $H^*(\mathrm{Fl}_n)$,
$\sigma_u\sigma_v=\sum_w c^w_{u,v}\sigma_w$; they vanish unless
$\ell(u)+\ell(v)=\ell(w)$.  The relevant duality is
$c^{\wo}_{u,v}=\delta_{v,\wo u}$ (verified below), so ``$\wo v$'' in Corollary~4
means $\wo\cdot v$.

\begin{theorem}[\texttt{math/0202090}, Corollary 4]\label{thm:cor4}
For all $u\le w$ in $\Sn$ and every type $\alpha$ with $|\alpha|=\ell(w)-\ell(u)$,
\[
  I_\alpha(u,w)\;=\;\sum_{v\,:\,\ell(v)=\ell(w)-\ell(u)} c^w_{u,v}\,
  I_\alpha(\wo v,\wo).
\]
\end{theorem}

\subsection{The conventions are validated, not assumed}

Reading a convention out of locators is where this kind of work goes wrong, so
the conventions above were \emph{tested before use}.  Specialising Theorem~2 of
\texttt{math/0202090} to $w=\wo$ predicts
\begin{equation}\label{eq:thm2}
  \mathfrak{S}_u(x)\;=\;\sum_{\gamma\in\Gam(u,\wo)} x^{\,\delta-\type(\gamma)},
  \qquad \delta=(n-1,n-2,\dots,1,0),
\end{equation}
with $\mathfrak{S}_u$ the Schubert polynomial.  Computing the right side from the
chain definitions above and the left side independently by divided differences
($\mathfrak{S}_{\wo}=x_1^{n-1}\cdots x_{n-1}$, $\mathfrak{S}_{ws_i}=\partial_i\mathfrak{S}_w$):

\begin{center}
\begin{tabular}{lrrr}
\toprule
 & permutations & agreeing & chains enumerated\\
\midrule
$n=2$ & 2 & 2 & 2\\
$n=3$ & 6 & 6 & 7\\
$n=4$ & 24 & 24 & 41\\
$n=5$ & 120 & 120 & 393\\
\midrule
total & 152 & \textbf{152} & 443\\
\bottomrule
\end{tabular}
\end{center}

\noindent
This is a genuine cross-mechanism check: the chain side uses only the labelled
Bruhat order, the polynomial side only divided differences.  An earlier reading
in which the monomial was $x^{\type(\gamma)}$ rather than $x^{\delta-\type(\gamma)}$
failed this test at once ($104$ of $120$ permutations disagreeing at $n=5$),
which is how the correct convention was found.  Everything downstream rests on
the validated version.

\section{Existence is trivial}

\begin{definition}\label{def:cor4map}
Fix $n$ and $u\le w$ in $\Sn$, $m=\ell(w)-\ell(u)$, and let
$T=\bigsqcup_{v:\ell(v)=m}\Gam(\wo v,\wo)$, a disjoint union with components
indexed by $v$.  A \emph{Corollary-4 map} for $(u,w)$ is a function
$f:\Gam(u,w)\to T$ such that
\begin{enumerate}
\item[(i)] $f$ is \emph{type-preserving}: $\type(f(\gamma))=\type(\gamma)$; and
\item[(ii)] for every $v$ with $\ell(v)=m$ and every $C\in\Gam(\wo v,\wo)$,
  $\;|f^{-1}(C)|=c^w_{u,v}$.
\end{enumerate}
This is exactly the object the remark after Theorem~2 asks for.
\end{definition}

\begin{theorem}\label{thm:A}
For every $n$ and every $u\le w$ in $\Sn$, Corollary-4 maps for $(u,w)$ exist,
and their number is
\[
  N(u,w)\;=\;\prod_{\alpha}\;
  \frac{I_\alpha(u,w)!}{\displaystyle\prod_{v}\bigl(c^w_{u,v}!\bigr)^{I_\alpha(\wo v,\wo)}},
\]
the products over all types $\alpha$ with $|\alpha|=m$ and all $v$ with
$\ell(v)=m$.  Moreover the existence of a Corollary-4 map is \emph{equivalent} to
Theorem~\ref{thm:cor4}.
\end{theorem}

\begin{proof}
A chain has exactly one type, so $\Gam(u,w)=\bigsqcup_\alpha\Gam_\alpha(u,w)$ and
likewise the type-$\alpha$ part of the $v$-component of $T$ is
$\Gam_\alpha(\wo v,\wo)$.  By (i), $f$ is therefore the disjoint union over
$\alpha$ of maps
\[
  f_\alpha:\;\Gam_\alpha(u,w)\;\longrightarrow\;
  P_\alpha:=\bigsqcup_{v}\Gam_\alpha(\wo v,\wo),
\]
and conditions for distinct $\alpha$ constrain disjoint parts of the data, hence
are independent.  Fix $\alpha$.  Specifying $f_\alpha$ is the same as specifying
its fibres, i.e.\ a family of pairwise disjoint subsets
$\{B_C\}_{C\in P_\alpha}$ of $\Gam_\alpha(u,w)$ with union
$\Gam_\alpha(u,w)$; condition (ii) says $|B_C|=c^w_{u,v}$ whenever
$C\in\Gam_\alpha(\wo v,\wo)$.  (For $v$ with $c^w_{u,v}=0$ every such $B_C$ is
empty, so the $v$-component receives nothing, as it must.)  Such a family of
prescribed sizes exists if and only if the sizes sum to the size of the set,
\[
  \sum_{C\in P_\alpha}|B_C|
  \;=\;\sum_{v} c^w_{u,v}\,I_\alpha(\wo v,\wo)
  \;\overset{!}{=}\;I_\alpha(u,w),
\]
which is precisely Theorem~\ref{thm:cor4} for this $\alpha$; and in that case the
number of such families is the multinomial coefficient
$I_\alpha(u,w)!\big/\prod_{C\in P_\alpha}|B_C|!
 = I_\alpha(u,w)!\big/\prod_v (c^w_{u,v}!)^{I_\alpha(\wo v,\wo)}$.
Types with $I_\alpha(u,w)=0$ have all blocks empty and contribute a factor $1$.
Multiplying over $\alpha$ gives $N(u,w)$, and $N(u,w)\ge1$ exactly when
Theorem~\ref{thm:cor4} holds for every $\alpha$.
\end{proof}

\begin{corollary}\label{cor:A}
The existence of a Corollary-4 map carries no information beyond
Theorem~\ref{thm:cor4}, and in particular yields no combinatorial
interpretation of $c^w_{u,v}$.
\end{corollary}

\begin{proof}
Immediate from Theorem~\ref{thm:A}: the two statements are equivalent.  The
structural reason is worth naming.  In Definition~\ref{def:cor4map} the numbers
$c^w_{u,v}$ appear as an \emph{input} --- they are the prescribed block sizes ---
and a combinatorial interpretation would have to produce them as an
\emph{output}, as the cardinality of a set defined without reference to them.
Theorem~\ref{thm:A} says the fibre condition is a bookkeeping constraint on a set
partition, satisfiable by fiat once the totals agree.
\end{proof}

\begin{remark}
This answers the literal target ``construct, or prove obstructed, a
type-preserving map $\Gam(u,w)\to\bigsqcup_v\Gam(\wo v,\wo)$ whose fibre over the
$v$-component has cardinality $c^w_{u,v}$'': it is constructed, for every $n$,
and the construction shows the target was mis-posed.  Lenart and Sottile's
sentence is not wrong --- they say such a map \emph{would give} a combinatorial
interpretation, meaning a described one.  The content of the 2002 problem lies
entirely in the word ``combinatorial'', i.e.\ in canonicity, and not in
existence.
\end{remark}

\section{For $n\le5$ the fibre condition is vacuous}

\begin{theorem}\label{thm:B}
\begin{enumerate}
\item For $2\le n\le 5$, every Schubert structure constant of $\mathrm{Fl}_n$
  satisfies $c^w_{u,v}\in\{0,1\}$.  \emph{(Exhaustive computation; two
  mechanisms.)}
\item For $n=6$, with $u=v=(1,3,5,2,4,6)$ and $w=(2,4,6,1,3,5)$ one has
  $c^w_{u,v}=2$.
\item Consequently the least $n$ for which some fibre of some Corollary-4 map has
  cardinality $\ge2$ is $n=6$.  For $n\le5$ every Corollary-4 map is a
  type-preserving \emph{bijection}
  $\Gam(u,w)\to\bigsqcup_{v:c^w_{u,v}=1}\Gam(\wo v,\wo)$, and
  $N(u,w)=\prod_\alpha I_\alpha(u,w)!$.
\end{enumerate}
\end{theorem}

\begin{proof}
(1) was computed exhaustively over all triples $(u,v,w)\in\Sn^3$ with
$\ell(u)+\ell(v)=\ell(w)$, for $n=3,4,5$, by expanding $\mathfrak{S}_u\mathfrak{S}_v$
and extracting the constant term of $\partial_w$; the maximum value attained was
$1$ in each case, with $0$ triples having $c\ge2$.  The computation was
cross-checked against Monk's rule
$\sigma_u\sigma_{s_r}=\sum_{a\le r<b,\;\ell(ut_{ab})=\ell(u)+1}\sigma_{ut_{ab}}$
on $7624$ comparisons ($852$ nonzero) with $0$ disagreements, and a negative
control in which $r$ was replaced by a different index produced $104$
disagreements at $n=4$, so the cross-check is not vacuous.

(2) $u$ and $w$ are the Grassmannian permutations of $\lambda=(2,1)$ and
$\lambda=(3,2,1)$ for $\mathrm{Gr}(3,6)$, of lengths $3$ and $6$.  Computing
directly in $\mathfrak{S}_6$ by divided differences gives $c^w_{u,u}=2$.  As an
independent check on a different mechanism entirely, enumerating
Littlewood--Richardson skew tableaux gives
$c^{(3,2,1)}_{(2,1),(2,1)}=2$, consistent with the classical expansion
$s_{21}^2=s_{42}+s_{411}+s_{33}+2s_{321}+s_{3111}+s_{222}+s_{2211}$, which the
same enumerator reproduces term for term.

(3) is immediate from (1), (2) and Theorem~\ref{thm:A}: if all $c^w_{u,v}\in\{0,1\}$
then every nonempty fibre is a singleton, the map is injective on each type, and
the multinomial collapses to $\prod_\alpha I_\alpha(u,w)!$.
\end{proof}

\begin{remark}[the plan for this session could not have worked]
The computation was scoped to ``$n=3$ and $n=4$, exhaustively''.
Theorem~\ref{thm:B} shows that no amount of computing in $\mathfrak{S}_3$ or
$\mathfrak{S}_4$ --- or $\mathfrak{S}_5$ --- can exhibit a single fibre of size
$\ge2$.  In that range the fibre-cardinality constraint, which is the entire
content of Definition~\ref{def:cor4map}(ii), is \emph{identically} the statement
that $f$ is a bijection.
\end{remark}

\section{Non-vacuity accounting}

Before searching, the number of \emph{falsifiable} instances was counted, under
the criterion that was specified in advance: triples $(u,w,\alpha)$ with
$I_\alpha(u,w)>1$ \emph{and} at least two distinct $v$ contributing.

\begin{center}
\begin{tabular}{lrrr}
\toprule
 & instances of Thm~\ref{thm:cor4} checked & disagreements & falsifiable\\
\midrule
$n=2$ & 3 & 0 & 0\\
$n=3$ & 23 & 0 & \textbf{0}\\
$n=4$ & 413 & 0 & \textbf{7}\\
\midrule
total & 439 & \textbf{0} & 7\\
\bottomrule
\end{tabular}
\end{center}

So Theorem~\ref{thm:cor4} is verified, with $0$ failures of the per-type identity
and $0$ failures of the coarser total-cardinality invariant
$|\Gam(u,w)|=\sum_v c^w_{u,v}|\Gam(\wo v,\wo)|$.  But $n=3$ has
\emph{no falsifiable instance at all}: every type-$\alpha$ chain set there has at
most one element, so a green sweep at $n=3$ carries exactly zero information.

Worse, the pre-specified criterion is itself too weak.  All seven $n=4$ instances
have the same shape: $I_\alpha(u,w)=2$, with exactly two contributing $v$, each
with $c^w_{u,v}=1$ and $I_\alpha(\wo v,\wo)=1$.  Both fibres are singletons;
there are $2!=2$ Corollary-4 maps and nothing in Definition~\ref{def:cor4map} to
prefer either.  The criterion with content is \emph{some $c^w_{u,v}\ge2$}, and by
Theorem~\ref{thm:B} that first occurs at $n=6$.

Explicitly, over all pairs $u\le w$ the number of Corollary-4 maps is:
$N(u,w)=1$ for all $23$ pairs at $n=3$; and at $n=4$, $N=1$ for $334$ pairs and
$N=2$ for the $7$ pairs above.  Every constructed map was re-verified against
both clauses of Definition~\ref{def:cor4map} independently of its construction
($23/23$ and $341/341$ reported \texttt{OK}).

\section{The uniqueness clause does not apply}

\begin{proposition}\label{prop:C}
The bijectivization-uniqueness clause of \texttt{2609.18502} \S5 ---
\emph{``A bijectivization always exists\ldots{} It is unique only when one of
$A,B$ is a singleton''} (l.~3497, read at source) --- does not apply to
Theorem~\ref{thm:cor4}, except on the locus $I_\alpha(u,w)=1$ where
$N(u,w)=1$ holds trivially.
\end{proposition}

\begin{proof}
Name the object on each side.  In \texttt{2609.18502}, a bijectivization of an
identity $\sum_{a\in A}\mathrm{wt}(a)=\sum_{b\in B}\mathrm{wt}(b)$ is a coupling
of the two sides: a joint measure on $A\times B$ with the prescribed marginals,
presented by stochastic kernels $\mathsf{p}^{\mathrm{fwd}},
\mathsf{p}^{\mathrm{bwd}}$.  The paper fixes the scope of $A$ and $B$ in its own
sentence: \emph{``$A$ and $B$ always mean the supports of the two sums:
configurations of weight zero are discarded, and all stochasticity and uniqueness
statements are imposed on the positive-weight support only.''}  So the hypothesis
is $|A|=1$ or $|B|=1$ as \emph{sets of configurations}.

For Theorem~\ref{thm:cor4} at a fixed type $\alpha$, the supports are
$A=\Gam_\alpha(u,w)$ and $B=\bigsqcup_v\{1,\dots,c^w_{u,v}\}\times
\Gam_\alpha(\wo v,\wo)$, with all weights equal to $1$.  Hence
$|A|=|B|=I_\alpha(u,w)$, and the hypothesis holds iff $I_\alpha(u,w)=1$.  That is
the degenerate locus, and by Theorem~\ref{thm:A} it is exactly where
$N=1$ anyway.

Two further mismatches hold even on that locus.  First, the object the clause
makes unique is a \emph{coupling}, not a bijection; the paper's own application
is explicit about this --- after collapsing one side to ``a \emph{single} term''
it concludes that the pair is \emph{``drawn with probability proportional to its
weight''} (l.~4229--4233).  A coupling is a bijection only in the zero-one
extreme, which for $|A|=1$ forces $|B|=1$ as well.  Second, the framework
contains no type-preservation constraint, so even a unique coupling would not be
a Corollary-4 map in the sense of Definition~\ref{def:cor4map}(i).

A quantitative replacement for the clause: with all weights $1$ and
$|A|=|B|=N$, the couplings of the two sides form $N\cdot\mathcal{B}_N$, $N$ times
the Birkhoff polytope of doubly stochastic $N\times N$ matrices, of dimension
$(N-1)^2$, whose vertices are precisely the bijections.  Uniqueness holds iff
$N=1$.
\end{proof}

\begin{remark}
The error the brief for this session made --- and it is the reason this phase was
scheduled --- was to read ``Corollary 4's left-hand side is a single
$I_\alpha(u,w)$'' as the singleton hypothesis.  A single \emph{term} of an
algebraic identity is not a singleton \emph{support}: $I_\alpha(u,w)$ is the
cardinality of $A$, not evidence that $|A|=1$.  The two readings share the word
and nothing else.
\end{remark}

\section{A negative result on the natural candidate}

Two structural facts make the chains unexpectedly rigid.  Both were verified
exhaustively for $n=3,4,5$ ($23$, $420$ and $17507$ chains):
\begin{enumerate}
\item A chain is determined by its source together with its label sequence:
  \emph{$0$ collisions}.
\item Because labels increase lexicographically, the first coordinates increase
  weakly, so the type $\alpha$ determines the first-coordinate \emph{sequence}
  exactly: $0$ violations.
\end{enumerate}
Hence a chain in $\Gam_\alpha(u,w)$ is faithfully encoded by its
second-coordinate word $b(\gamma)=(b_1,\dots,b_m)\in[n]^m$, and \emph{both sides
of Theorem~\ref{thm:cor4} are sets of words of the same length in the same
alphabet}.  This makes one candidate canonical map available for free: the
identity on words.  It is refuted.

\begin{proposition}
$\{b(\gamma):\gamma\in\Gam_\alpha(u,w)\}$ and
$\bigsqcup_v\{b(\gamma):\gamma\in\Gam_\alpha(\wo v,\wo)\}^{\,\oplus c^w_{u,v}}$
are in general distinct multisets of words: they differ in $7$ of $23$ instances
at $n=3$ and in $280$ of $413$ at $n=4$.  (The smallest witness: $n=3$,
$u=\mathrm{id}$, $w=(1,3,2)$, $\alpha=(0,1)$, left word $(2)$, right word $(1)$.)
\end{proposition}

No relabelling of the alphabet repairs this: at $n=3$ the failures demand
$2\mapsto1$ and $1\mapsto2$ at $m=1$, while at $m=2$ the instance
$u=\mathrm{id}$, $w=(2,3,1)$, $\alpha=(1,1)$ demands $(1,1)\mapsto(2,1)$, which
is not induced by any map of letters.

\section{Verification and controls}

\begin{center}
\begin{tabular}{lll}
\toprule
instrument & result & control\\
\midrule
chain apparatus vs.\ divided differences, \eqref{eq:thm2}
  & $152/152$ perms, $443$ chains
  & wrong monomial fails $104/120$\\
$c^w_{u,v}$: divided differences vs.\ Monk
  & $7624$ comparisons, $0$ disagree
  & wrong $r$: $104$ disagree\\
$c=2$ at $n=6$: $\mathfrak{S}_6$ vs.\ LR tableaux
  & both $2$
  & LR enumerator $9/9$ known values\\
Corollary 4
  & $439$ instances, $0$ disagree
  & total-cardinality invariant $0$ fail\\
constructed maps
  & $364/364$ re-verified \texttt{OK}
  & see below\\
\bottomrule
\end{tabular}
\end{center}

The three controls planted on the search, one per claim the apparatus makes:

\begin{description}
\item[C1, wrong type (target type permuted).]  Fired on $9/23$ at $n=3$ and
  $175/341$ at $n=4$.  The silences are accounted for exactly, with quantifiers:
  $4$ and $128$ pairs have no chains at all (not applicable); $8$ and $28$ have
  a shift-invariant type vector (the perturbation is the identity); and $2$ and
  $10$ have $\alpha\mapsto I_\alpha(u,w)$ \emph{constant} ($\equiv1$) on a
  shift-closed support, so permuting types permutes the support and preserves
  every count.  This last class is not a missed bug and not a vacuous control
  but a \emph{vacuous instance}: on such a pair \emph{no} type-permutation
  control can fire, and there are $10$ of $19$ such pairs at $n=3$ and $38$ of
  $213$ at $n=4$.
\item[C1b, wrong type of wrong degree (non-dodgeable).]  Because
  $|\alpha|=\ell(w)-\ell(u)$ is forced, perturbing the degree cannot be absorbed.
  Fired on $19/19$ and $213/213$ applicable pairs.
\item[C2, perturbed $c^w_{u,v}$ (one value incremented by $1$).]  Fired on
  $19/19$ and $213/213$ applicable pairs; $0$ silences.
\item[C3, chain enumerator off its domain.]  $\ell(w)<\ell(u)$: $0$ of $13$ and
  $0$ of $235$ pairs returned a chain.  $u\not\le w$ in Bruhat order with
  $\ell(w)>\ell(u)$: $0$ of $46$ at $n=4$ returned a chain.  Positive control
  ($u\le w$ must give chains): $0$ empty of $19$ and $213$.
\end{description}

\section{What is actually open}

Theorem~\ref{thm:A} removes existence from the problem and
Theorem~\ref{thm:B} removes $n\le5$ from the search space.  What remains is the
question Lenart and Sottile were actually pointing at, which should be stated
with the quantifier in the right place.

\begin{question}\label{q:canonical}
Is there a rule --- a map defined uniformly in $(u,w,\alpha)$ by a combinatorial
operation on chains, not referring to $c^w_{u,v}$ --- which for every $u\le w$
produces a Corollary-4 map?  Equivalently: can the $N(u,w)$ candidates of
Theorem~\ref{thm:A} be cut to one by a coherence condition, so that
$c^w_{u,v}$ becomes the \emph{output} (a fibre cardinality of a described map)
rather than the input?
\end{question}

Three things learned here constrain the search for such a rule.
\begin{enumerate}
\item The natural ambient set is the set of second-coordinate words: both sides
  of Theorem~\ref{thm:cor4} embed in $[n]^m$ faithfully, and the first
  coordinates are determined by $\alpha$ and carry no freedom.  A rule should be
  an operation on words.
\item It is not the identity on words (\S6), and not induced by any relabelling
  of letters.
\item The first instances where the fibre condition says anything at all are in
  $\mathfrak{S}_6$, where $c=2$ appears via $\mathrm{Gr}(3,6)\subset
  \mathrm{Fl}_6$.  Any candidate rule must be tested there; $\mathfrak{S}_4$
  cannot distinguish a correct rule from an arbitrary one, because at $n\le5$
  every Corollary-4 map is a bijection and $N(u,w)=\prod_\alpha I_\alpha(u,w)!$
  counts arbitrary relabellings.
\end{enumerate}

\section*{Gaps}

\begin{enumerate}
\item Theorem~\ref{thm:B}(1) is a finite exhaustive verification for
  $n\le5$, not a proof for all $n\le5$ by argument.  It is cross-checked by two
  mechanisms with a firing negative control, but it is a computation.  (The
  general statement ``$c^w_{u,v}\le1$ for $n\le5$'' presumably has a proof; I
  did not find one and did not look for one.)
\item Theorem~\ref{thm:cor4} itself is verified here only for $n\le4$ ($439$
  instances), and is quoted from \texttt{math/0202090} for general $n$.  I did
  not reprove it.
\item Question~\ref{q:canonical} is untouched.  Nothing in this note bears on
  whether a canonical rule exists; what is established is only that the
  existence formulation cannot be the question, and that $n\le5$ cannot be the
  test.
\item No computation was done at $n=6$ beyond the single structure constant of
  Theorem~\ref{thm:B}(2).  In particular the chain sets $\Gam_\alpha(u,w)$ in
  $\mathfrak{S}_6$ for the $\mathrm{Gr}(3,6)$ permutations were \emph{not}
  enumerated, so the first genuinely informative instance of
  Definition~\ref{def:cor4map} has not been examined.  That is the next step.
\end{enumerate}

\section*{Apparatus}

\texttt{proofs/code-1006c1-ls-cor4/}: \texttt{ls.py} (labelled Bruhat order,
increasing chains, types), \texttt{schub.py} (Schubert polynomials by divided
differences, structure constants), \texttt{lr.py} (independent
Littlewood--Richardson enumerator), \texttt{cor4.py} (Corollary 4 and the
non-vacuity gate), \texttt{search.py} (construction, independent re-verification,
map count).  Pure Python throughout, deliberately not Sage, so that the
ground truth runs on a different mechanism from \texttt{sage-combinat}.

\end{document}
